% Options for packages loaded elsewhere
% Options for packages loaded elsewhere
\PassOptionsToPackage{unicode}{hyperref}
\PassOptionsToPackage{hyphens}{url}
\PassOptionsToPackage{dvipsnames,svgnames,x11names}{xcolor}
%
\documentclass[
  spanish,
  11pt,
  a4paper,
]{article}
\usepackage{xcolor}
\usepackage[margin=2.5cm]{geometry}
\usepackage{amsmath,amssymb}
\setcounter{secnumdepth}{5}
\usepackage{iftex}
\ifPDFTeX
  \usepackage[T1]{fontenc}
  \usepackage[utf8]{inputenc}
  \usepackage{textcomp} % provide euro and other symbols
\else % if luatex or xetex
  \usepackage{unicode-math} % this also loads fontspec
  \defaultfontfeatures{Scale=MatchLowercase}
  \defaultfontfeatures[\rmfamily]{Ligatures=TeX,Scale=1}
\fi
\usepackage{lmodern}
\ifPDFTeX\else
  % xetex/luatex font selection
\fi
% Use upquote if available, for straight quotes in verbatim environments
\IfFileExists{upquote.sty}{\usepackage{upquote}}{}
\IfFileExists{microtype.sty}{% use microtype if available
  \usepackage[]{microtype}
  \UseMicrotypeSet[protrusion]{basicmath} % disable protrusion for tt fonts
}{}
\makeatletter
\@ifundefined{KOMAClassName}{% if non-KOMA class
  \IfFileExists{parskip.sty}{%
    \usepackage{parskip}
  }{% else
    \setlength{\parindent}{0pt}
    \setlength{\parskip}{6pt plus 2pt minus 1pt}}
}{% if KOMA class
  \KOMAoptions{parskip=half}}
\makeatother
% Make \paragraph and \subparagraph free-standing
\makeatletter
\ifx\paragraph\undefined\else
  \let\oldparagraph\paragraph
  \renewcommand{\paragraph}{
    \@ifstar
      \xxxParagraphStar
      \xxxParagraphNoStar
  }
  \newcommand{\xxxParagraphStar}[1]{\oldparagraph*{#1}\mbox{}}
  \newcommand{\xxxParagraphNoStar}[1]{\oldparagraph{#1}\mbox{}}
\fi
\ifx\subparagraph\undefined\else
  \let\oldsubparagraph\subparagraph
  \renewcommand{\subparagraph}{
    \@ifstar
      \xxxSubParagraphStar
      \xxxSubParagraphNoStar
  }
  \newcommand{\xxxSubParagraphStar}[1]{\oldsubparagraph*{#1}\mbox{}}
  \newcommand{\xxxSubParagraphNoStar}[1]{\oldsubparagraph{#1}\mbox{}}
\fi
\makeatother


\usepackage{longtable,booktabs,array}
\usepackage{calc} % for calculating minipage widths
% Correct order of tables after \paragraph or \subparagraph
\usepackage{etoolbox}
\makeatletter
\patchcmd\longtable{\par}{\if@noskipsec\mbox{}\fi\par}{}{}
\makeatother
% Allow footnotes in longtable head/foot
\IfFileExists{footnotehyper.sty}{\usepackage{footnotehyper}}{\usepackage{footnote}}
\makesavenoteenv{longtable}
\usepackage{graphicx}
\makeatletter
\newsavebox\pandoc@box
\newcommand*\pandocbounded[1]{% scales image to fit in text height/width
  \sbox\pandoc@box{#1}%
  \Gscale@div\@tempa{\textheight}{\dimexpr\ht\pandoc@box+\dp\pandoc@box\relax}%
  \Gscale@div\@tempb{\linewidth}{\wd\pandoc@box}%
  \ifdim\@tempb\p@<\@tempa\p@\let\@tempa\@tempb\fi% select the smaller of both
  \ifdim\@tempa\p@<\p@\scalebox{\@tempa}{\usebox\pandoc@box}%
  \else\usebox{\pandoc@box}%
  \fi%
}
% Set default figure placement to htbp
\def\fps@figure{htbp}
\makeatother



\ifLuaTeX
\usepackage[bidi=basic]{babel}
\else
\usepackage[bidi=default]{babel}
\fi
% get rid of language-specific shorthands (see #6817):
\let\LanguageShortHands\languageshorthands
\def\languageshorthands#1{}


\setlength{\emergencystretch}{3em} % prevent overfull lines

\providecommand{\tightlist}{%
  \setlength{\itemsep}{0pt}\setlength{\parskip}{0pt}}



 


\makeatletter
\@ifpackageloaded{tcolorbox}{}{\usepackage[skins,breakable]{tcolorbox}}
\@ifpackageloaded{fontawesome5}{}{\usepackage{fontawesome5}}
\definecolor{quarto-callout-color}{HTML}{909090}
\definecolor{quarto-callout-note-color}{HTML}{0758E5}
\definecolor{quarto-callout-important-color}{HTML}{CC1914}
\definecolor{quarto-callout-warning-color}{HTML}{EB9113}
\definecolor{quarto-callout-tip-color}{HTML}{00A047}
\definecolor{quarto-callout-caution-color}{HTML}{FC5300}
\definecolor{quarto-callout-color-frame}{HTML}{acacac}
\definecolor{quarto-callout-note-color-frame}{HTML}{4582ec}
\definecolor{quarto-callout-important-color-frame}{HTML}{d9534f}
\definecolor{quarto-callout-warning-color-frame}{HTML}{f0ad4e}
\definecolor{quarto-callout-tip-color-frame}{HTML}{02b875}
\definecolor{quarto-callout-caution-color-frame}{HTML}{fd7e14}
\makeatother
\makeatletter
\@ifpackageloaded{caption}{}{\usepackage{caption}}
\AtBeginDocument{%
\ifdefined\contentsname
  \renewcommand*\contentsname{Tabla de contenidos}
\else
  \newcommand\contentsname{Tabla de contenidos}
\fi
\ifdefined\listfigurename
  \renewcommand*\listfigurename{Listado de Figuras}
\else
  \newcommand\listfigurename{Listado de Figuras}
\fi
\ifdefined\listtablename
  \renewcommand*\listtablename{Listado de Tablas}
\else
  \newcommand\listtablename{Listado de Tablas}
\fi
\ifdefined\figurename
  \renewcommand*\figurename{Figura}
\else
  \newcommand\figurename{Figura}
\fi
\ifdefined\tablename
  \renewcommand*\tablename{Tabla}
\else
  \newcommand\tablename{Tabla}
\fi
}
\@ifpackageloaded{float}{}{\usepackage{float}}
\floatstyle{ruled}
\@ifundefined{c@chapter}{\newfloat{codelisting}{h}{lop}}{\newfloat{codelisting}{h}{lop}[chapter]}
\floatname{codelisting}{Listado}
\newcommand*\listoflistings{\listof{codelisting}{Listado de Listados}}
\makeatother
\makeatletter
\makeatother
\makeatletter
\@ifpackageloaded{caption}{}{\usepackage{caption}}
\@ifpackageloaded{subcaption}{}{\usepackage{subcaption}}
\makeatother
\usepackage{bookmark}
\IfFileExists{xurl.sty}{\usepackage{xurl}}{} % add URL line breaks if available
\urlstyle{same}
\hypersetup{
  pdftitle={Fuerzas Concurrentes en el Plano},
  pdfauthor={Cátedra: Ing. Vaello / Ing. Ronconi},
  pdflang={es},
  colorlinks=true,
  linkcolor={blue},
  filecolor={Maroon},
  citecolor={Blue},
  urlcolor={Blue},
  pdfcreator={LaTeX via pandoc}}


\title{Fuerzas Concurrentes en el Plano}
\usepackage{etoolbox}
\makeatletter
\providecommand{\subtitle}[1]{% add subtitle to \maketitle
  \apptocmd{\@title}{\par {\large #1 \par}}{}{}
}
\makeatother
\subtitle{TP N°1 --- Estabilidad 1 · 2° año Ingeniería Mecanica}
\author{Cátedra: Ing. Vaello / Ing. Ronconi}
\date{2026-04-09}
\begin{document}
\maketitle

\renewcommand*\contentsname{Tabla de contenidos}
{
\hypersetup{linkcolor=}
\setcounter{tocdepth}{3}
\tableofcontents
}

\section{Introducción: el concepto de
Resultante}\label{introducciuxf3n-el-concepto-de-resultante}

En aplicaciones reales --- un nudo de cercha, el anclaje de un puente,
la unión de una viga --- las fuerzas no actúan de a una. El nudo
experimenta la acción simultánea de múltiples fuerzas. La
\textbf{Resultante} es la fuerza única que produce el mismo efecto
mecánico que todo ese sistema.

\begin{tcolorbox}[enhanced jigsaw, bottomrule=.15mm, bottomtitle=1mm, titlerule=0mm, leftrule=.75mm, rightrule=.15mm, title=\textcolor{quarto-callout-note-color}{\faInfo}\hspace{0.5em}{Nota}, coltitle=black, colframe=quarto-callout-note-color-frame, left=2mm, opacitybacktitle=0.6, opacityback=0, toptitle=1mm, toprule=.15mm, colbacktitle=quarto-callout-note-color!10!white, colback=white, breakable, arc=.35mm]

El nudo no discrimina métodos analíticos ni gráficos. El nudo siente un
solo tirón: la Resultante.

\end{tcolorbox}

\textbf{Pregunta de apertura:} ¿Cambia el módulo de la Resultante si se
miden los ángulos desde el eje X o desde el eje Y?\\
→ No.~La Resultante es invariante respecto al sistema de referencia. Los
ejercicios 4 y 5 demuestran esto explícitamente.

\section{Ejercicios 1, 2 y 3 --- tres métodos sobre la misma
figura}\label{ejercicios-1-2-y-3-tres-muxe9todos-sobre-la-misma-figura}

La Figura 1 del TP presenta un sistema de 6 fuerzas concurrentes
(\(F_1\) a \(F_6\)). Cada uno de los primeros tres ejercicios trabaja
con un subconjunto distinto y aplica un método diferente para hallar la
Resultante.

\begin{longtable}[]{@{}lll@{}}
\toprule\noalign{}
Ejercicio & Subconjunto & Método \\
\midrule\noalign{}
\endhead
\bottomrule\noalign{}
\endlastfoot
1 & \(\{F_1, F_2, F_3, F_4, F_5\}\) & Proyecciones analíticas \\
2 & \(\{F_2, F_3, F_4, F_5, F_6\}\) & Momentos (Teorema de Varignon) \\
3 & \(\{F_1, F_3, F_4, F_5, F_6\}\) & Método mixto \\
\end{longtable}

Los tres ejercicios incluyen verificación gráfica mediante el polígono
de fuerzas.

\subsection{Ejercicio 1 --- Proyecciones
analíticas}\label{ejercicio-1-proyecciones-analuxedticas}

Las componentes cartesianas de la Resultante se obtienen sumando las
proyecciones de cada fuerza sobre los ejes X e Y:

\[R_x = \sum_{i=1}^{n} F_i \cos\theta_i \qquad R_y = \sum_{i=1}^{n} F_i \sin\theta_i\]

El módulo y la dirección resultan:

\[R = \sqrt{R_x^2 + R_y^2} \qquad \theta_R = \arctan\!\left(\frac{R_y}{R_x}\right)\]

\begin{tcolorbox}[enhanced jigsaw, bottomrule=.15mm, bottomtitle=1mm, titlerule=0mm, leftrule=.75mm, rightrule=.15mm, title=\textcolor{quarto-callout-warning-color}{\faExclamationTriangle}\hspace{0.5em}{Advertencia}, coltitle=black, colframe=quarto-callout-warning-color-frame, left=2mm, opacitybacktitle=0.6, opacityback=0, toptitle=1mm, toprule=.15mm, colbacktitle=quarto-callout-warning-color!10!white, colback=white, breakable, arc=.35mm]

\textbf{Condición implícita:} Pitágoras aplica únicamente cuando los
ejes son \textbf{ortogonales}. En el ejercicio 5 esa condición deja de
cumplirse.

\end{tcolorbox}

\subsection{Ejercicio 2 --- Momentos y Teorema de
Varignon}\label{ejercicio-2-momentos-y-teorema-de-varignon}

El momento de una fuerza respecto a un punto O es:

\[M_O(\mathbf{F}) = F \cdot d\]

donde \(d\) es el brazo de palanca (distancia perpendicular de O a la
línea de acción de \(\mathbf{F}\)).

\textbf{Teorema de Varignon:}

\[M_O(\mathbf{R}) = \sum_{i=1}^{n} M_O(\mathbf{F}_i)\]

El ejercicio pide plantear \textbf{dos ecuaciones de momentos} respecto
a puntos que no pertenezcan a los ejes coordenados.

\begin{tcolorbox}[enhanced jigsaw, bottomrule=.15mm, bottomtitle=1mm, titlerule=0mm, leftrule=.75mm, rightrule=.15mm, title=\textcolor{quarto-callout-tip-color}{\faLightbulb}\hspace{0.5em}{Tip}, coltitle=black, colframe=quarto-callout-tip-color-frame, left=2mm, opacitybacktitle=0.6, opacityback=0, toptitle=1mm, toprule=.15mm, colbacktitle=quarto-callout-tip-color!10!white, colback=white, breakable, arc=.35mm]

Elegir puntos fuera de los ejes fuerza a calcular todos los brazos de
palanca, lo que consolida el procedimiento. Si el punto estuviera sobre
un eje, algunas fuerzas tendrían brazo cero y el ejercicio perdería
riqueza.

\end{tcolorbox}

\begin{tcolorbox}[enhanced jigsaw, bottomrule=.15mm, bottomtitle=1mm, titlerule=0mm, leftrule=.75mm, rightrule=.15mm, title=\textcolor{quarto-callout-note-color}{\faInfo}\hspace{0.5em}{Nota}, coltitle=black, colframe=quarto-callout-note-color-frame, left=2mm, opacitybacktitle=0.6, opacityback=0, toptitle=1mm, toprule=.15mm, colbacktitle=quarto-callout-note-color!10!white, colback=white, breakable, arc=.35mm]

El método de momentos también determina la \textbf{ubicación de la recta
de acción} de la Resultante, información que las proyecciones no proveen
directamente.

\end{tcolorbox}

\subsection{Ejercicio 3 --- Método
mixto}\label{ejercicio-3-muxe9todo-mixto}

Combina una ecuación de proyección con una de momento.

\section{Ejercicios 4 y 5 --- invarianza de la
Resultante}\label{ejercicios-4-y-5-invarianza-de-la-resultante}

El sistema de fuerzas no cambia. Lo que cambia es el sistema de
referencia.

\subsection{Ejercicio 4 --- Rotación del sistema de
referencia}\label{ejercicio-4-rotaciuxf3n-del-sistema-de-referencia}

Se trabaja con la Resultante del ejercicio 1. La terna de ejes se rota
\(+45°\) en sentido antihorario, generando los ejes \(X' - Y'\).

El ejercicio pide calcular las nuevas componentes \(R_{x'}\) y
\(R_{y'}\), y demostrar que:

\[|\mathbf{R}| = \sqrt{R_{x'}^2 + R_{y'}^2}\]

coincide con el módulo obtenido en el ejercicio 1.

\begin{tcolorbox}[enhanced jigsaw, bottomrule=.15mm, bottomtitle=1mm, titlerule=0mm, leftrule=.75mm, rightrule=.15mm, title=\textcolor{quarto-callout-tip-color}{\faLightbulb}\hspace{0.5em}{Tip}, coltitle=black, colframe=quarto-callout-tip-color-frame, left=2mm, opacitybacktitle=0.6, opacityback=0, toptitle=1mm, toprule=.15mm, colbacktitle=quarto-callout-tip-color!10!white, colback=white, breakable, arc=.35mm]

Si el módulo cambia al rotar la terna, el error está en cómo se
calcularon los ángulos locales de cada fuerza respecto a los nuevos
ejes.

\end{tcolorbox}

\subsection{Ejercicio 5 --- Sistema de ejes
oblicuos}\label{ejercicio-5-sistema-de-ejes-oblicuos}

Se trabaja con el sistema del ejercicio 2. Los ejes \(X'' - Y''\) están
rotados \(+40°\) y \(+60°\) respectivamente en sentido \textbf{horario}
respecto a la terna original. El ángulo entre ambos ejes no es 90°.

\begin{tcolorbox}[enhanced jigsaw, bottomrule=.15mm, bottomtitle=1mm, titlerule=0mm, leftrule=.75mm, rightrule=.15mm, title=\textcolor{quarto-callout-warning-color}{\faExclamationTriangle}\hspace{0.5em}{Advertencia}, coltitle=black, colframe=quarto-callout-warning-color-frame, left=2mm, opacitybacktitle=0.6, opacityback=0, toptitle=1mm, toprule=.15mm, colbacktitle=quarto-callout-warning-color!10!white, colback=white, breakable, arc=.35mm]

\textbf{Prohibido Pitágoras.} Cuando los ejes no son perpendiculares, el
módulo de la Resultante se obtiene por la \textbf{ley de cosenos}:

\[R^2 = R_{x''}^2 + R_{y''}^2 + 2\,R_{x''}\,R_{y''}\cos\alpha\]

donde \(\alpha\) es el ángulo entre los ejes. El ejercicio pide
\textbf{justificar explícitamente} por qué Pitágoras no aplica.

\end{tcolorbox}

La verificación es gráfica mediante el \textbf{método del paralelogramo}
con los ejes oblicuos.

\section{Ejercicio 6 --- Descomposición y par de
transporte}\label{ejercicio-6-descomposiciuxf3n-y-par-de-transporte}

Se trabaja con la Resultante del ejercicio 3.

\subsection{Parte a) --- Descomposición en direcciones
oblicuas}\label{parte-a-descomposiciuxf3n-en-direcciones-oblicuas}

La Resultante se descompone en dos direcciones no ortogonales: dirección
\(a\) (a 30°) y dirección \(b\) (a 140°).

\begin{tcolorbox}[enhanced jigsaw, bottomrule=.15mm, bottomtitle=1mm, titlerule=0mm, leftrule=.75mm, rightrule=.15mm, title=\textcolor{quarto-callout-note-color}{\faInfo}\hspace{0.5em}{Nota}, coltitle=black, colframe=quarto-callout-note-color-frame, left=2mm, opacitybacktitle=0.6, opacityback=0, toptitle=1mm, toprule=.15mm, colbacktitle=quarto-callout-note-color!10!white, colback=white, breakable, arc=.35mm]

Como las direcciones no son perpendiculares entre sí, no se puede
aplicar la descomposición cartesiana estándar. El procedimiento es el
\textbf{método del paralelogramo} o la ley de senos.

\end{tcolorbox}

\subsection{Parte b) --- Traslación y par de
transporte}\label{parte-b-traslaciuxf3n-y-par-de-transporte}

Trasladar el punto de aplicación de una fuerza no es gratis. El sistema
original se mantiene equivalente solo si se agrega un \textbf{par de
transporte}:

\[M_{\text{par}} = \mathbf{R} \times d\]

El ejercicio pide trasladar la Resultante a los puntos \(A(-4,\,1)\) m y
\(B(2,\,3)\) m, determinando el par de transporte en cada caso.

La fuerza es idéntica en A y en B. Los pares son distintos porque las
distancias a la recta de acción original son distintas.

\section{Ejercicio 7 --- Equilibrio: el
semáforo}\label{ejercicio-7-equilibrio-el-semuxe1foro}

Un semáforo de peso \(P = 0{,}125\,\text{kN}\) cuelga de un cable
vertical unido a otros dos cables fijos a un soporte superior, que
forman ángulos de \(37°\) y \(53°\) con la horizontal.

La condición de equilibrio es:

\[\sum \mathbf{F}_i = \mathbf{0} \quad \Leftrightarrow \quad \sum F_x = 0 \quad \text{y} \quad \sum F_y = 0\]

Se pide determinar la tracción en los tres cables.

\textbf{Interpretación gráfica:} el polígono de fuerzas cierra. La
última fuerza regresa exactamente al punto de partida del diagrama;
cualquier desvío indica error.

\begin{tcolorbox}[enhanced jigsaw, bottomrule=.15mm, bottomtitle=1mm, titlerule=0mm, leftrule=.75mm, rightrule=.15mm, title=\textcolor{quarto-callout-warning-color}{\faExclamationTriangle}\hspace{0.5em}{Advertencia}, coltitle=black, colframe=quarto-callout-warning-color-frame, left=2mm, opacitybacktitle=0.6, opacityback=0, toptitle=1mm, toprule=.15mm, colbacktitle=quarto-callout-warning-color!10!white, colback=white, breakable, arc=.35mm]

El enunciado da el peso en \textbf{kN}. Trabajar en kN durante todo el
ejercicio.

\end{tcolorbox}

\section{Ejercicio 8 --- Equilibrante de cada
sistema}\label{ejercicio-8-equilibrante-de-cada-sistema}

\begin{longtable}[]{@{}
  >{\raggedright\arraybackslash}p{(\linewidth - 2\tabcolsep) * \real{0.5000}}
  >{\raggedright\arraybackslash}p{(\linewidth - 2\tabcolsep) * \real{0.5000}}@{}}
\toprule\noalign{}
\begin{minipage}[b]{\linewidth}\raggedright
Concepto
\end{minipage} & \begin{minipage}[b]{\linewidth}\raggedright
Definición
\end{minipage} \\
\midrule\noalign{}
\endhead
\bottomrule\noalign{}
\endlastfoot
\textbf{Resultante} \(\mathbf{R}\) & Fuerza única equivalente al
sistema. Lo que \emph{le pasa} al nudo. \\
\textbf{Equilibrante} \(\mathbf{E}\) & Fuerza necesaria para anular la
Resultante. Lo que \emph{se aplica} para que el nudo no se mueva. \\
\end{longtable}

\[\mathbf{E} = -\mathbf{R}\]

El ejercicio pide determinar la equilibrante para los tres sistemas de
los ejercicios 1, 2 y 3, usando en cada caso el mismo método que se
empleó para hallar la Resultante:

\begin{longtable}[]{@{}
  >{\raggedright\arraybackslash}p{(\linewidth - 2\tabcolsep) * \real{0.5000}}
  >{\raggedright\arraybackslash}p{(\linewidth - 2\tabcolsep) * \real{0.5000}}@{}}
\toprule\noalign{}
\begin{minipage}[b]{\linewidth}\raggedright
Sistema
\end{minipage} & \begin{minipage}[b]{\linewidth}\raggedright
Método
\end{minipage} \\
\midrule\noalign{}
\endhead
\bottomrule\noalign{}
\endlastfoot
\(\mathbf{E}_1\) (del Ej. 1) & Analítico + verificación por polígono
cerrado \\
\(\mathbf{E}_2\) (del Ej. 2) & Momentos, verificar \(\sum M = 0\)
respecto a un punto arbitrario \\
\(\mathbf{E}_3\) (del Ej. 3) & Método mixto \\
\end{longtable}

\begin{tcolorbox}[enhanced jigsaw, bottomrule=.15mm, bottomtitle=1mm, titlerule=0mm, leftrule=.75mm, rightrule=.15mm, title=\textcolor{quarto-callout-warning-color}{\faExclamationTriangle}\hspace{0.5em}{Advertencia}, coltitle=black, colframe=quarto-callout-warning-color-frame, left=2mm, opacitybacktitle=0.6, opacityback=0, toptitle=1mm, toprule=.15mm, colbacktitle=quarto-callout-warning-color!10!white, colback=white, breakable, arc=.35mm]

\textbf{Error frecuente:} la equilibrante actúa en el \textbf{mismo
punto de aplicación} que la Resultante, con sentido contrario. No es
simplemente cambiar el signo del vector en el gráfico sin verificar el
punto.

\end{tcolorbox}

\section{Ejercicio 9 --- Variación del momento y Teorema de
Varignon}\label{ejercicio-9-variaciuxf3n-del-momento-y-teorema-de-varignon}

Usando la Resultante del ejercicio 1, el ejercicio estudia cómo varía el
efecto de giro según el punto respecto al cual se mide.

Se pide:

\begin{enumerate}
\def\labelenumi{\alph{enumi})}
\tightlist
\item
  Calcular el momento de \(\mathbf{R}\) respecto a un punto genérico
  \(P(0,\,y)\) sobre el eje de ordenadas.
\item
  Determinar la coordenada \(y_P\) para la cual el momento sea de
  \(20\,\text{kN·m}\) en sentido horario.
\item
  Verificar el resultado aplicando el Teorema de Varignon con las
  fuerzas individuales del subconjunto.
\end{enumerate}

\begin{tcolorbox}[enhanced jigsaw, bottomrule=.15mm, bottomtitle=1mm, titlerule=0mm, leftrule=.75mm, rightrule=.15mm, title=\textcolor{quarto-callout-note-color}{\faInfo}\hspace{0.5em}{Nota}, coltitle=black, colframe=quarto-callout-note-color-frame, left=2mm, opacitybacktitle=0.6, opacityback=0, toptitle=1mm, toprule=.15mm, colbacktitle=quarto-callout-note-color!10!white, colback=white, breakable, arc=.35mm]

El momento de una fuerza no es una propiedad fija: depende del punto
respecto al cual se mide. Varignon permite calcularlo descomponiendo la
fuerza en sus componentes, lo que suele simplificar el cálculo de brazos
de palanca.

\end{tcolorbox}

\section{Síntesis}\label{suxedntesis}

\begin{enumerate}
\def\labelenumi{\arabic{enumi}.}
\tightlist
\item
  La Resultante es invariante respecto al sistema de referencia
  (demostrado en ej. 4 y 5).
\item
  Proyecciones, momentos y método mixto son herramientas
  complementarias; el método mixto es la herramienta de auditoría.
\item
  En sistemas oblicuos, Pitágoras no aplica --- usar ley de cosenos (ej.
  5).
\item
  Descomponer en direcciones no ortogonales requiere paralelogramo o ley
  de senos (ej. 6a).
\item
  Trasladar fuerzas introduce momentos adicionales --- par de transporte
  (ej. 6b).
\item
  Equilibrio equivale a Resultante nula y polígono vectorial cerrado
  (ej. 7).
\item
  Equilibrante \(= -\mathbf{R}\): mismo módulo y línea de acción,
  sentido contrario, mismo punto de aplicación (ej. 8).
\item
  El momento varía con el punto de referencia; Varignon permite
  calcularlo por descomposición (ej. 9).
\end{enumerate}




\end{document}
